% !TeX root = ../../Book1.tex
% 纲要标题：### 20.3 Kummer 理论
% 纲要位置：计划纲要.md，第 720 行。

\section{Kummer 理论}
\label{b1:sec:2003}

方程 \(X^n-a=0\) 看起来只涉及开方，但它的全部根还相差单位根。只有先确定基域中是否包含这些单位根，根式扩张与循环群之间的关系才足够直接。\term[Kummer-lilun]{Kummer 理论}[Kummer theory]研究这类根式与交换 Galois 群的联系。本节固定代数闭包 \(\overline K\)，所有扩张都取在其中；除另有说明外，\(n\geq1\)、\(\operatorname{char}K\nmid n\)，并假设 \(K\) 含有全部 \(n\) 次单位根。

% 纲要内容（仅作写作计划，尚非教材正文）：
%
% * 特征不整除 n 且基域含所需单位根的假设
% * 先以 Hilbert 第 90 定理证明循环扩张的根式生成，作为第21.2节的直接准备
% * 单根式扩张次数等于被开方数的商类阶；根式参数选择与 Kummer 配对的良定义
% * 有限 Kummer 对应的双侧单射、阶比较与同时特征基重建；素数次扩张由参数直线分类
% * 在本节定义 \(\mu_n\) 及本原单位根；循环扩张的根式生成不依赖一般 Kummer 对应，也不调用后续分圆多项式不可约性
%

\subsection{特征与单位根假设}
\label{b1:subsec:2003:01}

\begin{definition}\label{b1:def:roots-of-unity}
设 \(K\) 为域，\(\overline K\) 为 \(K\) 的代数闭包，\(n\geq1\) 为整数。如果 \(\zeta\in\overline K^\times\) 满足 \(\zeta^n=1\)，则称 \(\zeta\) 为 \(n\) 次\term[danwei-gen]{单位根}[root of unity]。全体 \(n\) 次单位根组成乘法群
\[
\mu_n\coloneqq\{\,\zeta\in\overline K^\times\mid\zeta^n=1\,\}.
\]
如果 \(\zeta\in\mu_n\) 的乘法阶恰为 \(n\)，则称 \(\zeta\) 为\term[benyuan-danwei-gen]{本原单位根}[primitive root of unity]。
\end{definition}
\symidx{\(\mu_n\)：全体 \(n\) 次单位根组成的群}

在 \(\operatorname{char}K\nmid n\) 的假设下，\(X^n-1\) 与导数 \(nX^{n-1}\) 没有公共根，所以 \(\mu_n\) 有 \(n\) 个元素。\cref{b1:lem:finite-multiplicative-cyclic}说明域的乘法群的有限子群是循环群；\(\mu_n\) 因而是阶为 \(n\) 的循环群，它的任一生成元 \(\zeta\) 都有阶 \(n\)，即为本原 \(n\) 次单位根，且 \(\mu_n=\langle\zeta\rangle\)。假设 \(\mu_n\subseteq K\) 就等价于 \(K\) 含有这样一个 \(\zeta\)。

\ProseParagraphBreak
对 \(a\in K^\times\)，特征条件使 \((X^n-a)'=nX^{n-1}\) 与 \(X^n-a\) 没有公共根，保证多项式可分；若去掉它，在特征 \(p\) 中就有 \(X^p-a=(X-\alpha)^p\)，其中 \(\alpha^p=a\)，只有一个重根，并且 \(X^p-1=(X-1)^p\) 表明不存在阶为 \(p\) 的单位根。单位根条件则保证：已有一个根 \(\alpha\) 后，其余根 \(\zeta\alpha\)（\(\zeta\in\mu_n\)）全部位于 \(K(\alpha)\) 中。若去掉它，\(\mathbb Q(\sqrt[3]{2})/\mathbb Q\) 就是可分而不正规的根式扩张。这两个假设分别排除了重根与遗漏其他共轭根的障碍。

\subsection{循环扩张与根式生成}
\label{b1:subsec:2003:04}

\begin{proposition}\label{b1:prop:single-radical-cyclic}
设 \(K\) 为域，\(n\geq1\) 为整数，\(\operatorname{char}K\nmid n\)，且 \(K\) 包含全部 \(n\) 次单位根。固定 \(K\) 的代数闭包 \(\overline K\)，取 \(a\in K^\times\) 及满足 \(\alpha^n=a\) 的 \(\alpha\in\overline K\)。则 \(K(\alpha)/K\) 是循环 Galois 扩张，其次数整除 \(n\)。映射
\[
\operatorname{Gal}\bigl(K(\alpha)/K\bigr)\longrightarrow\mu_n,
\qquad\sigma\longmapsto\frac{\sigma(\alpha)}{\alpha}
\]
是单同态。
\end{proposition}
\begin{proof}
取 \(\mu_n\) 的生成元 \(\zeta\)。\(X^n-a\) 的根恰为 \(\zeta^j\alpha\)，全在 \(K(\alpha)\) 中；因 \(a\neq0\) 且 \(n\neq0\) 在 \(K\) 中成立，它没有重根。所以 \(K(\alpha)/K\) 是一个可分多项式的分裂域，为有限 Galois 扩张。对其任一自同构 \(\sigma\)，
\[
\sigma(\alpha)^n=\sigma(\alpha^n)=\sigma(a)=a=\alpha^n,
\]
故 \((\sigma(\alpha)/\alpha)^n=1\)，比值属于 \(\mu_n\)。又因 \(\mu_n\subseteq K\)，每个自同构都固定这些比值，所以
\[
\frac{\sigma\tau(\alpha)}{\alpha}
=\sigma\Bigl(\frac{\tau(\alpha)}{\alpha}\Bigr)\frac{\sigma(\alpha)}{\alpha}
=\frac{\tau(\alpha)}{\alpha}\frac{\sigma(\alpha)}{\alpha}.
\]
目标群交换，故上式就是普通群同态的乘法公式。比值为 \(1\) 时，自同构固定生成元 \(\alpha\)，因而为恒等映射，这证明了单射性。Galois 群由此同构于 \(\mu_n\) 的一个子群，所以为循环群，群阶整除 \(|\mu_n|=n\)。由于扩张已经证明为 Galois，\([K(\alpha):K]=|\operatorname{Gal}(K(\alpha)/K)|\)，扩张次数也整除 \(n\)。
\end{proof}

这个同态计算再次用到了单位根条件。在一个更大的分裂域中，若自同构不固定单位根，则 \(\sigma(\tau(\alpha)/\alpha)\) 不能直接换成 \(\tau(\alpha)/\alpha\)，仍须保留自同构对该比值的作用；那正是上一节一次余圈的乘法公式。

\ProseParagraphBreak
根式扩张的次数还可以从被开方数在乘法商群中的阶直接读出。

\begin{corollary}\label{b1:cor:kummer-radical-degree}
设 \(K\) 为域，\(n\geq1\) 为整数，\(\operatorname{char}K\nmid n\)，且 \(K\) 包含全部 \(n\) 次单位根。固定 \(K\) 的代数闭包，取 \(a\in K^\times\)，并在该闭包中取满足 \(\alpha^n=a\) 的根 \(\alpha\)。记 \(K^{\times n}=\{\,c^n\mid c\in K^\times\,\}\)。则
\[
[K(\alpha):K]=\operatorname{ord}_{K^\times/K^{\times n}}([a]).
\]
\end{corollary}
\begin{proof}
记 \(L=K(\alpha)\)、\(d=[L:K]\)，并令 \(e\) 为 \([a]\) 的阶。因为 \([a]^n=[1]\)，\(e\) 是整除 \(n\) 的正整数。由\cref{b1:prop:single-radical-cyclic}，\(L/K\) 循环 Galois，其 Galois 群通过比值映射嵌入 \(\mu_n\)，像群的阶为 \(d\)。因此每个比值 \(\sigma(\alpha)/\alpha\) 的 \(d\) 次幂为一，故 \(\sigma(\alpha^d)=\alpha^d\) 对所有自同构成立。不动域为 \(K\)，所以 \(\alpha^d\in K^\times\)，进而 \(a^d=(\alpha^d)^n\in K^{\times n}\)，即 \(e\mid d\)。

反过来，\(a^e=c^n\) 对某个 \(c\in K^\times\) 成立，故 \((\alpha^e/c)^n=1\)。由 \(\mu_n\subseteq K\) 得 \(\alpha^e/c\in K\)，所以 \(\alpha^e\in K\)。于是 \(\alpha\) 满足 \(K\) 上的 \(e\) 次多项式 \(X^e-\alpha^e\)，其最小多项式次数 \(d\) 不超过 \(e\)。结合 \(e\mid d\) 与 \(d\le e\)，得到 \(d=e\)。
\end{proof}

反过来，循环扩张也可以用一个根式生成。

\begin{theorem}\label{b1:thm:kummer-cyclic}
设 \(L/K\) 是次数为 \(d\) 的循环扩张。如果 \(\operatorname{char}K\nmid d\)，并且 \(K\) 包含全部 \(d\) 次单位根 \(\mu_d\)，则存在 \(\alpha\in L^\times\) 使
\[
L=K(\alpha),\qquad\alpha^d\in K.
\]
\end{theorem}

为寻找这样的 \(\alpha\)，选 Galois 群的生成元 \(\sigma\) 和本原 \(d\) 次单位根 \(\zeta\)，要求 \(\sigma(\alpha)=\zeta\alpha\)，就是在线性代数意义下寻找 \(\sigma\) 的一个特征向量，特征值为 \(\zeta\)。这使 \(\sigma(\alpha^d)=\zeta^d\alpha^d=\alpha^d\)，又把构造化为 \(\alpha/\sigma(\alpha)=\zeta^{-1}\)，恰可调用乘法形式的 Hilbert 第 90 定理。

\begin{proof}
选取 \(\operatorname{Gal}(L/K)\) 的生成元 \(\sigma\)，以及本原 \(d\) 次单位根 \(\zeta\in K\)。因为 \(\zeta\) 属于基域，全部自同构都固定它，所以
\[
N_{L/K}(\zeta^{-1})=\prod_{j=0}^{d-1}\sigma^j(\zeta^{-1})
=\zeta^{-d}=1.
\]
\cref{b1:thm:hilbert90-multiplicative}给出非零 \(\alpha\) 满足 \(\alpha/\sigma(\alpha)=\zeta^{-1}\)，即 \(\sigma(\alpha)=\zeta\alpha\)。所以 \(\alpha^d\) 被 \(\sigma\) 固定，属于 \(K\)。又有 \(\sigma^j(\alpha)=\zeta^j\alpha\) 对 \(0\le j<d\) 成立。\(\sigma^j\) 固定最小多项式 \(m_{\alpha,K}\) 的系数，故这些像全部是它的根；因为 \(\zeta\) 的阶为 \(d\) 且 \(\alpha\ne0\)，它们两两不同。因此 \(\deg m_{\alpha,K}\ge d\)。另一方面，\(K(\alpha)\subseteq L\) 给出 \(\deg m_{\alpha,K}=[K(\alpha):K]\le[L:K]=d\)，故次数相等，\(K(\alpha)=L\)。
\end{proof}

\cref{b1:thm:kummer-cyclic}将用于\secref{b1:sec:2102}的特征零根式可解性判据。多个根式之间的关系还可由有限交换 Galois 群统一描述，下面据此建立 Kummer 配对。

\subsection{根式扩张与 Kummer 配对}
\label{b1:subsec:2003:02}

一个根式的自同构由比值 \(\sigma(\alpha)/\alpha\) 决定。加入多个根式时，我们希望同时记录这些比值，并判断哪些根式增加了扩张次数。若 \(a'=ab^n\)、\(b\in K^\times\)，而 \(\alpha^n=a\)，则 \((b\alpha)^n=a'\)，且 \(K(b\alpha)=K(\alpha)\)。因此应先把被开方数放入商群 \(K^\times/K^{\times n}\)，再用同一自同构在各个根上的比值来比较它们。

\ProseParagraphBreak
为同时处理若干根式，记 \(K^{\times n}=\{\,b^n\mid b\in K^\times\,\}\)。商群 \(K^\times/K^{\times n}\) 中每个元素的阶都整除 \(n\)。取其有限子群 \(B\)，选择生成元的代表 \(a_1,\ldots,a_r\)，并令
\[
L_B=K(\alpha_1,\ldots,\alpha_r),\qquad\alpha_i^n=a_i.
\]
换一个根只会乘以 \(\mu_n\subseteq K\) 中的元素；换一个代表只会再乘以 \(K^\times\) 中的元素。若 \([a]\in B\)，可写 \([a]=[a_1]^{e_1}\cdots[a_r]^{e_r}\)，其中 \(e_i\in\mathbb Z\)。这意味着对某个 \(c\in K^\times\) 有
\[
a=c^n a_1^{e_1}\cdots a_r^{e_r},\qquad
\bigl(c\alpha_1^{e_1}\cdots\alpha_r^{e_r}\bigr)^n=a.
\]
所有 \(\alpha_i\) 均非零，所以负指数也有意义；\(a\) 的任一 \(n\) 次根只与所列的根相差 \(K\) 中的单位根，故都属于 \(L_B\)。换一组生成元时，新生成元代表的根因而全部属于原来的 \(L_B\)，反过来原生成元的根也属于新域，所以两域相等。这说明 \(L_B\) 不依赖代表、根或生成列表的选择。它也是多项式 \(X^n-a_1,\ldots,X^n-a_r\) 的共同分裂域；这些多项式都可分，故 \(L_B/K\) 是有限 Galois 扩张。记 \(G_B=\operatorname{Gal}(L_B/K)\)。

\begin{definition}\label{b1:def:kummer-pairing}
设 \(K\) 为域，\(n\geq1\) 且 \(\operatorname{char}K\nmid n\)，在 \(K\) 的固定代数闭包 \(\overline K\) 中记 \(\mu_n=\{\zeta\in\overline K^\times\mid\zeta^n=1\}\)，并假设 \(\mu_n\subseteq K\)。取乘法商群 \(K^\times/K^{\times n}\) 的有限子群 \(B\)，其中 \(K^{\times n}=\{b^n\mid b\in K^\times\}\)；令 \(L_B\) 为在 \(\overline K\) 中给 \(B\) 的生成元代表添入 \(n\) 次根所得的域，并记 \(G_B=\operatorname{Gal}(L_B/K)\)。对 \(\sigma\in G_B\) 与 \(\xi=[a]\in B\)，取 \(a\in K^\times\) 及满足 \(\alpha^n=a\) 的 \(\alpha\in L_B\)，定义
\[
\langle\sigma,\xi\rangle\coloneqq\frac{\sigma(\alpha)}{\alpha},
\qquad \alpha^n=a.
\]
该值与代表 \(a\) 及根 \(\alpha\) 的选择无关。由此得到的映射 \(G_B\times B\rightarrow\mu_n\) 称为\term[Kummer-peidui]{Kummer 配对}[Kummer pairing]。
\end{definition}

配对的值与根和代表的选择无关。若把 \(\alpha\) 换成 \(\alpha'=\zeta\alpha\)，其中 \(\zeta\in\mu_n\subseteq K\)，则
\[
\frac{\sigma(\alpha')}{\alpha'}=
\frac{\zeta\sigma(\alpha)}{\zeta\alpha}=
\frac{\sigma(\alpha)}{\alpha}.
\]
若把代表换成 \(a'=ab^n\)，其中 \(b\in K^\times\)，可取根 \(b\alpha\)，同样有 \(\sigma(b\alpha)/(b\alpha)=\sigma(\alpha)/\alpha\)。该代表的其他根再由前一种计算处理，故两个选择都不影响配对值。

\ProseParagraphBreak
根的乘积是乘积的根，所以配对在第二个变量上保持乘法；第一个变量上的计算仍须使用 \(\sigma\) 固定 \(\mu_n\) 这一事实，正如\cref{b1:prop:single-radical-cyclic}证明中的比值公式。于是对 \(\sigma,\tau\in G_B\)、\(\xi,\eta\in B\)，有
\[
\begin{aligned}
\langle\sigma\tau,\xi\rangle&=
\langle\sigma,\xi\rangle\langle\tau,\xi\rangle,\\
\langle\sigma,\xi\eta\rangle&=
\langle\sigma,\xi\rangle\langle\sigma,\eta\rangle.
\end{aligned}
\]
因此固定任意一个变量，就得到另一个变量上的群同态。

\subsection{有限 Kummer 对应}
\label{b1:subsec:2003:03}

在 \(\operatorname{char}K\nmid n\) 的假设下，\(\mu_n\) 是阶为 \(n\) 的循环群。设 \(A\) 为有限交换群。对群同态 \(\chi,\psi:A\rightarrow\mu_n\)，\(\operatorname{Hom}(A,\mu_n)\) 的运算是逐点相乘：\((\chi\psi)(x)=\chi(x)\psi(x)\)，单位元是常值 \(1\) 的同态。\(A\) 的指数指使每个元素的该次幂等于 \(1\) 的最小正整数。若其指数整除 \(n\)，则
\begin{equation}\label{b1:eq:finite-character-count}
\bigl|\operatorname{Hom}(A,\mu_n)\bigr|=|A|.
\end{equation}
这里应用\cref{b1:thm:fga-structure}，把有限交换群 \(A\) 写成若干循环群的直积；因 \(A\) 的指数整除 \(n\)，每个循环因子的阶 \(d\) 也整除 \(n\)。从 \(C_d\) 到 \(\mu_n\) 的同态由生成元的像决定，这个像必须满足 \(\zeta^d=1\)。在阶为 \(n\) 的循环群中，满足这个条件的元素有 \(\gcd(d,n)\) 个；只有 \(d\mid n\) 时才恰有 \(d\) 个。当前每个循环因子均满足此条件，各因子的选择独立，相乘得到所需等式。这说明指数条件不能省去；单凭 \(A\) 有限交换，并不能保证有 \(|A|\) 个这样的同态。

\begin{theorem}[Kummer 对应]\label{b1:thm:kummer-correspondence}
设 \(K\) 为域，\(n\geq1\) 为整数，\(\operatorname{char}K\nmid n\)，并且 \(K\) 包含全部 \(n\) 次单位根 \(\mu_n\)。固定 \(K\) 的代数闭包 \(\overline K\)，记 \(K^{\times n}=\{c^n\mid c\in K^\times\}\)。以下两类对象一一对应：
\begin{enumerate}
\item \(K^\times/K^{\times n}\) 的有限子群 \(B\)；
\item \(\overline K\) 中的有限 Galois 扩张 \(L/K\)，其 Galois 群为指数整除 \(n\) 的交换群。
\end{enumerate}
对第一类中的有限子群 \(B\)，令 \(L_B\subseteq\overline K\) 为给 \(B\) 中各元素的代表添入 \(n\) 次根所生成的域。正向对应为 \(B\mapsto L_B\)，反向对应为
\[
B_L=\left(K^\times\cap L^{\times n}\right)/K^{\times n}.
\]
这里 \(L^{\times n}=\{\,\beta^n\mid\beta\in L^\times\,\}\)，所以分子收集 \(K\) 中那些在 \(L\) 内已经具有 \(n\) 次根的非零元素，再模掉它们在 \(K\) 中本来就是 \(n\) 次幂的平凡类。
Kummer 配对给出同构
\[
\operatorname{Gal}(L_B/K)\cong\operatorname{Hom}(B,\mu_n),
\qquad B\cong\operatorname{Hom}\bigl(\operatorname{Gal}(L_B/K),\mu_n\bigr),
\]
特别地，\([L_B:K]=|B|\)。子群包含对应于域包含。
\end{theorem}
\begin{proof}
从有限 \(B\) 出发，配对给出群同态
\[
G_B\longrightarrow\operatorname{Hom}(B,\mu_n),\qquad
\sigma\longmapsto\bigl(\xi\mapsto\langle\sigma,\xi\rangle\bigr).
\]
若 \(\sigma\) 在其核中，则全部选定根的比值 \(\sigma(\alpha_i)/\alpha_i\) 均为一。因此 \(\sigma\) 固定生成域的每个根 \(\alpha_i\)，也固定由它们和 \(K\) 生成的整个 \(L_B\)，故为恒等。这个映射因而单射。目标是逐点相乘的交换群，每个同态的 \(n\) 次幂都是常值一的同态，所以这个单射还说明 \(G_B\) 交换且指数整除 \(n\)。

在另一侧，配对给出 \(B\rightarrow\operatorname{Hom}(G_B,\mu_n)\)。若 \([a]\) 对每个 \(\sigma\in G_B\) 的配对值都是一，则其根 \(\alpha\) 被整个群固定。不动域为 \(K\)，所以 \(\alpha\in K\)，\(a=\alpha^n\) 在 \(K\) 中本来就是 \(n\) 次幂，即 \([a]=[1]\)。这也得到一个单射：任何非平凡被开方数类都能被某个自同构通过配对区别出来。

有限交换群 \(B\) 的指数整除 \(n\)，所以\eqref{b1:eq:finite-character-count}给出 \(|\operatorname{Hom}(B,\mu_n)|=|B|\)，第一个单射推出 \(|G_B|\le|B|\)。刚才已从这个单射得到 \(G_B\) 的交换性与指数条件，因而同一个计数公式也适用于 \(G_B\)，得到 \(|\operatorname{Hom}(G_B,\mu_n)|=|G_B|\)。第二个单射于是推出 \(|B|\le|G_B|\)。两边的有限阶相等，两个单射都为满射，故都是同构。又因 \(L_B/K\) 为 Galois，\([L_B:K]=|G_B|=|B|\)。

现在从第二类扩张 \(L/K\) 出发，记群为 \(G\)。要将它重建为一个根式扩张，只需找到一组 \(K\) 基 \(\beta_i\)，使每个 \(\beta_i^n\) 都在 \(K\) 中。把各 \(\sigma\in G\) 视为 \(K\) 向量空间 \(L\) 上的线性算子，寻找它们的共同特征向量，就能使这些 \(n\) 次幂被整个群固定。由于群的指数整除 \(n\)，各算子满足 \(\sigma^n=1\)；\(X^n-1\) 在 \(K\) 中分裂且无重根，可以用插值多项式把向量分解到其特征空间中。对每个 \(\zeta\in\mu_n\)，令
\[
P_\zeta(X)=\prod_{\substack{\eta\in\mu_n\\\eta\neq\zeta}}
\frac{X-\eta}{\zeta-\eta}.
\]
这就是\cref{b1:thm:lagrange-interpolation}在互异插值点 \(\mu_n\) 上的基多项式，满足 \(P_\zeta(\eta)=1\) 当 \(\eta=\zeta\)，否则为零。因而 \(\sum_\zeta P_\zeta(X)-1\) 是次数小于 \(n\) 却有 \(n\) 个不同根的多项式，只能为零。又有
\[
(X-\zeta)P_\zeta(X)=
\frac{X^n-1}{\prod_{\eta\in\mu_n,\,\eta\ne\zeta}(\zeta-\eta)},
\]
因为分子恰是全部 \(n\) 个一次因子的乘积，分母是非零的基域常数。代入算子 \(\sigma\)，利用 \(\sigma^n=1\)，得到 \((\sigma-\zeta I)P_\zeta(\sigma)=0\)。所以
\[
v=\sum_{\zeta\in\mu_n}P_\zeta(\sigma)v
\]
的每一项都属于 \(\sigma\) 的 \(\zeta\) 特征空间。\(P_\zeta(\sigma)\) 在该空间上为恒等，在其他特征空间上为零；若这些空间中向量之和为零，分别作用各 \(P_\zeta(\sigma)\) 就得到每项都为零。因此这个分解是直和，\(\sigma\) 可对角化。

要得到整个 \(G\) 的同时特征基，还需要群交换。若 \(\sigma(v)=\zeta v\)，则对任意 \(\tau\in G\)，
\[
\sigma(\tau(v))=\tau(\sigma(v))=\tau(\zeta v)=\zeta\tau(v),
\]
因为两算子交换且 \(\zeta\in K\)。因此一个算子的每个特征空间都被其余算子保持，可以在其中对下一算子再作上述直和分解。群 \(G\) 有限，依次处理其全部元素后，便得到 \(L\) 的公共特征空间直和分解；在每个非零公共特征空间中选一组基，合在一起得到同时特征向量基 \(\beta_1,\ldots,\beta_m\)。这里指数条件用于单个算子的分解，交换性用于使这些分解相容。

写 \(\sigma(\beta_i)=\lambda_{\sigma,i}\beta_i\)，其中 \(\lambda_{\sigma,i}\in K\)。因为 \(\beta_i\ne0\)，
\[
\beta_i=\sigma^n(\beta_i)=\lambda_{\sigma,i}^n\beta_i
\]
给出 \(\lambda_{\sigma,i}^n=1\)，即 \(\lambda_{\sigma,i}=\sigma(\beta_i)/\beta_i\in\mu_n\)。于是 \(\sigma(\beta_i^n)=\beta_i^n\) 对所有 \(\sigma\in G\) 成立，故 \(\beta_i^n\in K^\times\)。这些 \(n\) 次幂的类生成一个有限子群 \(B_0\)：各生成元的幂指数都可约为 \(0,\ldots,n-1\)，所以其元素数至多为 \(n^m\)。构造 \(L_{B_0}\) 时可取 \(\beta_i\) 作为这些类的根，因此 \(L_{B_0}\subseteq L\)。另一方面，这个子域含有 \(K\) 和全部 \(\beta_i\)，故含有它们的所有 \(K\) 线性组合，即整个 \(L\)。两边包含给出 \(L_{B_0}=L\)。

\(B_0\) 是利用一次基的选择构造出来的参数群；还须确认它已经包含 \(L\) 中所有被开方数类，即等于由扩张本身定义的 \(B_L\)。对 \(B_L\) 中任意元素，在 \(L\) 中选取其代表的 \(n\) 次根。以同一比值公式得到群同态 \(B_L\rightarrow\operatorname{Hom}(G,\mu_n)\)；若某个类与所有 \(\sigma\in G\) 的配对值均为一，其根便属于不动域 \(K\)，所以该映射单射。目标是有限群，因而 \(B_L\) 有限。由构造可知 \(B_0\subseteq B_L\)；而 \(B_L\) 中各类在 \(L\) 中有根，另选的根只差 \(K\) 中的单位根，因此 \(L_{B_L}\subseteq L\)。于是
\[
L=L_{B_0}\subseteq L_{B_L}\subseteq L.
\]
故 \(L_{B_L}=L\)。对两个有限子群分别应用已证的次数公式，得到 \(|B_L|=[L:K]=|B_0|\)。因为 \(B_0\subseteq B_L\)，有限群的阶相等迫使 \(B_L=B_0\)，所以选取的这组基虽然不唯一，所得到的参数群却是扩张内在确定的全部参数群。

反过来，从任意有限 \(B\) 出发，由各代表在 \(L_B\) 中具有根可知 \(B\subseteq B_{L_B}\)。刚才对第二类扩张的证明应用于 \(L_B/K\)，给出 \(B_{L_B}\) 有限且 \(L_{B_{L_B}}=L_B\)。因此 \(|B_{L_B}|=[L_B:K]=|B|\)，由包含与阶相等得 \(B_{L_B}=B\)。这证明了两个构造互逆。若 \(B\subseteq B'\)，则其代表的根也包含于 \(L_{B'}\)，从而 \(L_B\subseteq L_{B'}\)；反之，若 \(L_B\subseteq L_{B'}\)，一个在前者中已有 \(n\) 次根的基域元素在后者中也有根，故 \(B=B_{L_B}\subseteq B_{L_{B'}}=B'\)。
\end{proof}

这里参数子群 \(B\) 增大，所要添入的根也增多，所以 Kummer 对应中的子群包含与域包含同向。上一章的 Galois 对应则反向：自同构子群越大，共同不动元需要满足的条件越多，不动域越小。两种对应中的子群分别记录被开方数类与自同构，含义不同。

\ProseParagraphBreak
定理要求扩张有限，且 Galois 群的指数整除给定的 \(n\)。与之对应的 \(B\) 是有限子群，但不要求它在整个 \(K^\times/K^{\times n}\) 中具有有限指数。例如取 \(K=\mathbb Q\)、\(n=2\)，则 \(\mathbb Q^\times/\mathbb Q^{\times2}\) 是无限群，但 \(B=\{[1],[2],[3],[6]\}\) 有四个元素，给出 \(L_B=\mathbb Q(\sqrt2,\sqrt3)\)。独立性可由 \(2,3,6\) 都不是有理数的平方核验。

\ProseParagraphBreak
令 \(s\) 只改变 \(\sqrt2\) 的符号，\(t\) 只改变 \(\sqrt3\) 的符号。此时 \(G_B=\{1,s,t,st\}\)、\(\mu_2=\{1,-1\}\)，配对的全部取值为
\[
\begin{array}{c|rrrr}
\langle\sigma,\xi\rangle &[1]&[2]&[3]&[6]\\ \hline
1  & 1& 1& 1& 1\\
s  & 1&-1& 1&-1\\
t  & 1& 1&-1&-1\\
st & 1&-1&-1& 1
\end{array}
\]
每一行是 \(B\rightarrow\mu_2\) 的同态，每一列是 \(G_B\rightarrow\mu_2\) 的同态；各行两两不同，各列也两两不同，具体显示了配对的两侧非退化性。特别地，\([6]\) 的核为 \(\{1,st\}\)，所以 \(\mathbb Q(\sqrt6)=L_B^{\langle st\rangle}\)。这里 \(\{[1],[6]\}\subseteq B\) 给出子域包含 \(\mathbb Q(\sqrt6)\subseteq L_B\)，而子域与其固定子群之间的包含方向相反。

\ProseParagraphBreak
设 \(K\) 为域，\(\ell\) 为素数。乘法商群 \(V=K^\times/K^{\times\ell}\) 的每个元素的阶都整除 \(\ell\)，因而可以把 \(V\) 视为 \(\mathbb F_\ell\) 向量空间；\(V\) 也可能是平凡群。若用加法记号，其运算为
\[
[a]+[b]=[ab],\qquad \bar r\,[a]=[a^r],\qquad 0=[1]
\quad(a,b\in K^\times,\ r\in\mathbb Z).
\]
其中 \(\bar r\) 是整数 \(r\) 在 \(\mathbb F_\ell\) 中的像。因为 \([a]^\ell=[1]\)，若 \(r'\equiv r\pmod\ell\)，则 \([a^{r'}]=[a^r]\)，所以标量乘法只依赖 \(\bar r\)，不依赖其整数代表。分配律与标量乘法的结合律来自交换群的幂运算规则，例如 \((ab)^r=a^rb^r\)、\(a^{r+s}=a^ra^s\)。这里的非零向量是 \([a]\ne[1]\)，即 \(a\in K^\times\setminus K^{\times\ell}\)；只写 \(a\ne0\) 并不足以保证这一点。


\begin{corollary}\label{b1:cor:kummer-prime-classification}
设 \(K\) 为域，\(\ell\) 为不同于 \(\operatorname{char}K\) 的素数，且 \(K\) 包含全部 \(\ell\) 次单位根 \(\mu_\ell\)。在 \(K\) 的一个固定代数闭包中，\(K\) 的 \(\ell\) 次循环扩张由 \(K^\times/K^{\times\ell}\) 的一维 \(\mathbb F_\ell\) 子空间分类。
\end{corollary}
\begin{proof}
素数情形应用\cref{b1:thm:kummer-correspondence}：次数为 \(\ell\) 的域对应阶为 \(\ell\) 的子群。上述乘法商群的指数整除 \(\ell\)，它作为 \(\mathbb F_\ell\) 向量空间时，阶为 \(\ell\) 的子群正是一维子空间。
\end{proof}

因此，若 \(a\in K^\times\setminus K^{\times\ell}\)，同一直线上的 \(\ell-1\) 个非零向量 \([a]^r\)（\(1\le r<\ell\)）都生成同一个参数子群 \(\langle[a]\rangle\)，给出同一个非平凡 \(\ell\) 次扩张。具体地，若 \(\alpha^\ell=a\)、\(b\in K^\times\)，则 \(b\alpha^r\) 的 \(\ell\) 次幂为 \(a^r b^\ell\)。由于 \(r\) 与 \(\ell\) 互素，取整数 \(u,v\) 使 \(ru+\ell v=1\)，就有
\[
\alpha=(\alpha^r)^u a^v\in K(b\alpha^r).
\]
反向包含显然，所以 \(K(b\alpha^r)=K(\alpha)\)。以每个非零类分别标记扩张，会将这一条直线重复计算 \(\ell-1\) 次；一维子空间恰好把这些生成元变化归为同一个参数。这个区别与下一节的加法分类相对应。
